% -------------------------------------------------------
\section{Integration with the BX3 Framework}
\label{sec:bx3_integration}
% -------------------------------------------------------

The Recursive Spawning Protocol is a structural instantiation of the BX3 three-layer accountability architecture, not a module composed alongside it as a supplementary component. Every RSP mechanism maps onto exactly one BX3 layer, and this mapping is minimal and structurally necessary: no RSP guarantee can be achieved with fewer than all three layers, and no layer can subsume another's role without collapsing at least one correctness property. This section specifies the three mappings and establishes why each is exclusive to its layer.

% -------------------------------------------------------
\subsection{Purpose Layer: Accountability Anchor}
\label{sec:purpose_anchor}

The \purpose{} layer occupies two structurally singular roles in RSP, each constitutive of the protocol's correctness guarantees and each exclusive to \purpose{} by architectural definition.

\medskip
\noindent\textbf{Exclusive origin of spawn authorization.} At root provisioning, the \purpose{} layer defines and records the spawn authorization policy $P_{\mathrm{spawn}}$ in the \fact{} layer authorization ledger. $P_{\mathrm{spawn}}$ encodes two mandatory preconditions for any child spawn: (i) the child operating scope must be a strict descendant refinement of the parent scope ($R(n_{i+1}) \subset R(n_i)$), and (ii) the spawn must be operationally necessary --- the unmet condition must be unsatisfiable within the parent's existing scope. These conditions originate exclusively at the \purpose{} layer. The \bounds{} engine evaluates proposed spawns against $P_{\mathrm{spawn}}$; the \fact{} layer enforces it; neither originates it. No machine actor may issue a spawn warrant in the absence of a matching \purpose{} authorization record. The separation is architectural, not organizational: it holds regardless of privilege level, operational urgency, or system state.

\medskip
\noindent\textbf{Exclusive terminus of escalation.} When the chain reaches $D_{\mathrm{ESCALATE}}$, the protocol compulsory-escalates to the human anchor $h(n)$. The human anchor is non-collapsing: for all nodes $n_i$ in chain $C$, $h(n_i) = h(n_0)$. The anchor does not migrate down the chain as depth increases. The human issues one of three directives --- \texttt{ACK}, \texttt{RESOLVE}, or \texttt{REJECT} --- and no machine actor may issue any of these. The chain terminates in human judgment, never in autonomous resolution.

This is the structural expression of the upstream BX3 accountability guarantee: systems built on the BX3 Framework fail upward into human consciousness. If any machine actor could absorb an exception or issue a terminal directive, the accountability chain would terminate at that actor rather than at a human. The \purpose{} layer's exclusivity here is load-bearing: removing it causes the upstream accountability guarantee to collapse into a machine-actor loop.

\medskip
\noindent\textbf{Structural necessity.} The \purpose{} layer's two roles are jointly exclusive to it because both require the exercise of intent and the bearing of final accountability --- capabilities structurally reserved for the human principal. The \bounds{} engine cannot originate $P_{\mathrm{spawn}}$ because it operates within, not above, the operating range defined by the \purpose{} layer; it has no mandate to define what counts as authorized purpose. The \fact{} layer cannot serve as the escalation terminus because it is architecturally a deterministic enforcement layer, not a judgment layer: it applies rules, it does not decide.

% -------------------------------------------------------
\subsection{Bounds Engine: Depth Enforcement}
\label{sec:bounds_depth}

The \bounds{} engine provides constrained execution evaluation: it evaluates proposed actions against the operating range defined by the \purpose{} layer and enforces the depth constraint that prevents unbounded chain growth. Critically, the \bounds{} engine cannot exceed these constraints under any operational circumstance --- there is no privilege level, operational exigency, or system state under which it may override the depth bound.

\medskip
\noindent\textbf{Depth evaluation at spawn.} At each spawn event, the \bounds{} engine evaluates current chain depth $|C|$ against the \purpose{}-defined maximum $D_{\mathrm{max}}$, both anchored in the \fact{} layer authorization ledger. If $|C| \geq D_{\mathrm{max}}$, the \bounds{} engine does not execute the spawn. It transitions to \texttt{ESCALATING} state and initiates the Bailout Protocol, propagating the exception upward to $h(n)$. This enforcement is deterministic: the \bounds{} engine has no discretion to exceed $D_{\mathrm{max}}$ regardless of urgency, operational exigency, or the importance of the unmet condition. The depth constraint is a hard architectural boundary, not a heuristic.

\medskip
\noindent\textbf{Scope refinement evaluation.} The \bounds{} engine also evaluates the scope subset condition $R(n_{i+1}) \subset R(n_i)$ at each spawn event, determining whether a spawn request is justified by operational necessity within the \purpose{} layer's authorization scope. This evaluation is a precondition for \fact{} layer pre-commit hook invocation; it is not sufficient alone, because the scope refinement condition is independently enforced by the \fact{} layer pre-commit hook. This reflects the depth-limit insufficiency theorem: depth limits alone cannot prevent scope creep drift, so both depth enforcement and scope refinement enforcement are required as structurally distinct operations.

\medskip
\noindent\textbf{Fact Layer pre-commit reinforcement.} The \bounds{} engine's depth evaluation is structurally reinforced by the \fact{} layer pre-commit hook. Before any spawn operation is recorded, the \fact{} layer independently verifies: (a) $R(n_{i+1}) \subseteq R(n_i)$ (scope subset), and (b) $\mathrm{depth}(n_{i+1}) \leq D_{\mathrm{max}}$ (depth bound). The \fact{} layer rejects any spawn violating either condition. The \bounds{} engine cannot execute a rejected spawn; it may only transition to \texttt{ESCALATING}. The pre-commit hook makes the depth bound and scope refinement constraint structural invariants --- not policy conventions overridable in exigent circumstances.

% -------------------------------------------------------
\subsection{Fact Layer: Forensic Ledger}
\label{sec:fact_ledger}

The \fact{} layer provides deterministic enforcement and forensic auditability. It maintains the append-only ledger that makes the RSP chain auditable at any future time and enforces the write protocol constraints that make RSP structurally tamper-evident. The \fact{} layer is the structural foundation on which both the \purpose{} layer's authorization policy and the \bounds{} engine's depth constraints become runtime-enforceable, not merely documented.

\medskip
\noindent\textbf{Append-only forensic ledger.} The \fact{} layer records every spawn event, every state transition, every \purpose{} layer directive received, and every unresolved exception. Each entry is timestamped, attributed to the node that generated it, and hash-chained to the preceding entry. Hash-chaining ensures that no entry can be inserted, deleted, reordered, or altered after recording: any tampering breaks hash continuity and is detectable by any observer with ledger read access. The ledger is the single source of truth for all authorized chain state --- the runtime artifact that distinguishes RSP chains from conventional delegation chains, which exist only as forensic reconstructions attempted after the fact.

The forensic ledger closes the accountability loop: it records every directive issued by the human anchor and every resulting state transition, creating a complete, tamper-evident provenance chain from the root human principal to every active node in the system.

\medskip
\noindent\textbf{Immutable parent pointer.} Once $\alpha(n)$ is established at node provisioning, no actor --- including the \purpose{} layer after provisioning --- may modify it. The \fact{} layer write protocol rejects any attempt to overwrite $\alpha(n)$. Chain topology changes occur only through authorized spawn operations recorded in the ledger; there is no mechanism for retroactive pointer manipulation. The parent pointer is not merely access-controlled --- it is structurally immutable by protocol definition. This eliminates the re-parenting drift pathway by making retroactive chain rewriting structurally impossible regardless of requestor identity or privilege level.

\medskip
\noindent\textbf{Pre-commit hook.} The pre-commit hook is the enforcement mechanism that makes all other RSP constraints structural. At every spawn event, it verifies conditions (a) and (b) before the spawn is recorded. Rejected spawns are logged; the \bounds{} engine does not execute them; the protocol transitions to \texttt{ESCALATING} if the unmet condition persists. The pre-commit hook operates identically on all machine actors; there is no privileged actor capable of bypassing it. Even administrative compromise of a machine actor cannot produce a bypass, because the hook is enforced at the \fact{} layer write protocol level, not at the actor's discretion.

% -------------------------------------------------------
\subsection{Structural Minimality}
\label{sec:minimality}

The three-layer mapping is minimal and sufficient. Removing any layer causes the corresponding guarantee to fail (Table~\ref{tab:layer_removal}).

\begin{table}[h]
\begin{center}
\begin{tabular}{p{3.2cm} p{4.8cm}}
\toprule
\textbf{Removed layer} & \textbf{Guarantee that fails} \\
\midrule
\purpose{} & Exclusive human terminus eliminated; chain can terminate in machine actors; spawn policy becomes self-authorized by machine actors \\
\bounds{} & Constrained execution evaluation eliminated; no layer to evaluate spawn necessity or manage escalation transition \\
\fact{} & Immutable parent pointer eliminated; retroactive chain manipulation possible; scope refinement and depth bound become policy conventions, not structural invariants \\
\bottomrule
\end{tabular}
\caption{Layer removal consequences for RSP correctness properties.}
\label{tab:layer_removal}
\end{center}
\end{table}

The three-layer architecture is the minimal structure that makes RSP's correctness properties unconditionally guaranteed. No subset of layers achieves the same guarantees.

% -------------------------------------------------------
% -------------------------------------------------------
\section{Correctness Properties}
\label{sec:correctness}
% -------------------------------------------------------

This section establishes three correctness properties for the Recursive Spawning Protocol: drift prevention, chain integrity, and termination. Each property is stated formally, proved sketchedly, and connected to the BX3 Framework layers that enforce it. Together these properties constitute the RSP's overall guarantee: a spawn chain is an accountable, bounded, and auditable refinement process that terminates in human judgment.

% -------------------------------------------------------
\subsection{Drift Prevention}
\label{sec:drift_prevention}

\medskip
\noindent\textbf{Theorem (Drift Prevention).} Let $C = \langle n_0, n_1, \ldots, n_k \rangle$ be a Recursive Spawning chain rooted at $n_0$ with human anchor $h(n_0) = h(n_k)$. Then for all $i$ ($0 \leq i \leq k$), the operating scope $R(n_i)$ is a descendant refinement of the intent of $h(n_0)$. Equivalently: no node in $C$ can exhibit behavior outside the authorized intent of its human anchor.

\medskip
\noindent\textbf{Proof sketch (structural induction).} The proof proceeds by induction on chain depth, where each inductive step is enforced at runtime by the \fact{} layer pre-commit hook.

\medskip
\textit{Base case ($i=0$).} $R(n_0)$ is defined by $h(n_0)$ at provisioning. By the \purpose{} layer's definition of authorized operating scope, $R(n_0)$ is a descendant refinement of $h(n_0)$'s intent by construction. No drift exists at the root.

\medskip
\textit{Inductive step.} Assume $R(n_i)$ is a descendant refinement of $h(n_0)$'s intent for some $i$ with $0 \leq i < k$. For $n_i$ to spawn $n_{i+1}$, the \fact{} layer pre-commit hook verifies two conditions before the spawn is recorded:

\begin{enumerate}
    \item[(a)] $R(n_{i+1}) \subseteq R(n_i)$ --- the scope refinement constraint, required by $P_{\mathrm{spawn}}$; there is no authorized provision for lateral or upward scope expansion.
    \item[(b)] $\mathrm{depth}(n_{i+1}) \leq D_{\mathrm{max}}$ --- the depth constraint, also enforced by the pre-commit hook.
\end{enumerate}

The \fact{} layer rejects any spawn violating (a) or (b). No actor --- including the \bounds{} engine --- can execute a rejected spawn. Therefore $R(n_{i+1}) \subseteq R(n_i)$ holds by structural enforcement, not by policy convention.

By the inductive hypothesis, $R(n_i)$ is a descendant refinement of $h(n_0)$'s intent. Since $R(n_{i+1}) \subseteq R(n_i)$, transitivity gives that $R(n_{i+1})$ is also a descendant refinement of $h(n_0)$'s intent. Hence $n_{i+1}$ does not drift. $\square$

\medskip
The structural interpretation clarifies why this argument is immune to the failure modes that defeat policy-based approaches. Drift prevention does not rest on the reliability of any actor. It rests on the conjunction of three architectural facts, each necessary and jointly sufficient:

\begin{enumerate}
    \item[(1)] \textbf{Immutable parent pointer.} $\alpha(n)$ is established once at provisioning and is thereafter immutable under the \fact{} layer write protocol. Scope inheritance is traceable to a unique authorized lineage, not a manipulable graph. If the parent pointer were mutable, a node could re-parent to a different lineage and inherit a different scope, breaking the refinement chain.
    \item[(2)] \textbf{Forensic ledger.} Every scope assignment is recorded in the append-only, hash-chained ledger. The ledger makes the inductive proof empirically verifiable: each step of the refinement chain is a recorded, tamper-evident ledger entry. Without the ledger, the inductive argument is an assertion about recorded state; with it, the argument is a verifiable property of the historical record.
    \item[(3)] \textbf{Chain terminates at \purpose{} layer.} The exclusive human terminus closes the inductive argument at a finite bound. There is no machine-actor alternative terminus. If the chain could terminate in an autonomous actor, that actor could issue a terminal directive that escapes the refinement chain entirely.
\end{enumerate}

If any of these three conditions is absent, drift prevention is not structurally guaranteed. With all three, drift is architecturally impossible.

% -------------------------------------------------------
\subsection{Chain Integrity}
\label{sec:chain_integrity}

\medskip
\noindent\textbf{Theorem (Chain Integrity).} For any RSP chain $C = \langle n_0, n_1, \ldots, n_k \rangle$, the parent pointer sequence $\alpha(n_i) = n_{i-1}$ holds for all $1 \leq i \leq k$, and no node configuration in $C$ has been modified after provisioning.

\medskip
\noindent\textbf{Proof sketch.} The proof has two components: topological integrity and node immutability.

\medskip
\textit{Topological integrity.} Each spawn event $n_{i-1} \xrightarrow{\mathrm{spawn}} n_i$ establishes $\alpha(n_i) = n_{i-1}$ as an immutable \fact{} layer property at provisioning time. The \fact{} layer write protocol rejects any operation attempting to modify $\alpha(n_i)$ after provisioning. Since the chain is constructed exclusively through authorized spawn events and no parent pointer is modified after establishment, $\alpha(n_i) = n_{i-1}$ holds for all $i \geq 1$ by construction. The parent pointer is not merely access-controlled --- there is no privileged actor, including the \purpose{} layer after provisioning, capable of overwriting $\alpha(n)$. This eliminates the re-parenting drift pathway by making retroactive chain rewriting structurally impossible.

\medskip
\textit{Node immutability.} Each node $n_i$ is provisioned with a fixed configuration $(R(n_i), P_{\mathrm{spawn}}(n_i), h(n_i))$ established by the \purpose{} layer and recorded in the \fact{} layer authorization ledger. The \fact{} layer enforces immutability of these parameters after provisioning. The write protocol rejects any modification request from any actor --- including the \bounds{} engine and the \purpose{} layer after provisioning. The \purpose{} layer may issue a \texttt{RESOLVE} directive that redirects a node's behavior, but this directive is recorded as a new ledger entry (a state transition), not as an alteration of the provisioned configuration.

The forensic ledger's hash-chaining provides evidentially for chain integrity: any attempt to insert a falsified spawn event or reorder existing entries breaks hash continuity and is detectable by any observer with ledger read access. Chain integrity is not merely a property of the current chain state --- it is a property of the entire historical record, enforceable at any future time. $\square$

% -------------------------------------------------------
\subsection{Termination}
\label{sec:termination}

\medskip
\noindent\textbf{Theorem (Termination).} Every RSP chain $C = \langle n_0, n_1, \ldots, n_k \rangle$ terminates --- either at a resolution state or at the escalation threshold --- within a finite number of spawn steps bounded by $\max(D_{\mathrm{max}}, D_{\mathrm{ESCALATE}})$.

\medskip
\noindent\textbf{Proof sketch.} The proof establishes two lemmas and combines them.

\medskip
\textit{Lemma 1 (Depth-bounded growth).} The chain grows only through spawn events. By the \bounds{} engine depth enforcement and the \fact{} layer pre-commit hook, any spawn event producing $|C| \geq D_{\mathrm{max}}$ is rejected. Therefore the chain cannot exceed depth $D_{\mathrm{max}} - 1$ through autonomous spawn. Rejected spawns are logged; the \bounds{} engine transitions to \texttt{ESCALATING}. Autonomous chain growth is thus bounded by $D_{\mathrm{max}}$.

\medskip
\textit{Lemma 2 (Escalation-triggered halt).} If the chain reaches $D_{\mathrm{ESCALATE}}$ before resolving the exception condition, the protocol enters \texttt{ESCALATING} state and suspends all autonomous spawn activity. The \bounds{} engine enters quiescent hold, awaiting a directive from $h(n_k)$ --- the human anchor. No further spawn events occur until $h(n_k)$ issues \texttt{ACK}, \texttt{RESOLVE}, or \texttt{REJECT}. Since $h(n_k) = h(n_0)$ is a human principal, the directive arrives in finite time. \texttt{REJECT} terminates the chain definitively; \texttt{RESOLVE} redirects behavior and records a new ledger entry; \texttt{ACK} permits continued operation under \purpose{} layer authorization.

\medskip
\textit{Theorem conclusion.} Combining Lemma 1 and Lemma 2: the chain grows autonomously at most $D_{\mathrm{max}} - 1$ spawn steps. If it reaches $D_{\mathrm{ESCALATE}}$ before resolution, it halts and escalates. If it reaches $D_{\mathrm{max}}$ before resolution, the \fact{} layer rejects further spawn and triggers mandatory escalation. In all execution paths, the chain reaches a terminal state within $\max(D_{\mathrm{max}}, D_{\mathrm{ESCALATE}})$ spawn steps. Since both $D_{\mathrm{max}}$ and $D_{\mathrm{ESCALATE}}$ are finite integers defined at provisioning, the chain terminates in finite time. $\square$

The termination guarantee also precludes infinite spawn-escalate cycles. When $D_{\mathrm{max}}$ is reached, the \fact{} layer blocks further spawn and triggers mandatory escalation. The human anchor's \texttt{REJECT} directive terminates the chain. There is no mechanism by which the chain can circumvent this bound and re-enter autonomous spawn --- the \fact{} layer pre-commit hook is structural, not configurable at runtime by any machine actor.

% -------------------------------------------------------
\subsection{Collective Guarantee}
\label{sec:collective_guarantee}

The three correctness properties are mutually reinforcing and jointly enforced by the BX3 Framework's three-layer architecture.

\textbf{Drift prevention} ensures that scope refinement is the only authorized direction of chain growth. Without it, successive scope expansions could gradually migrate the chain's operating scope away from the human anchor's intent --- an autonomy runway within an otherwise bounded-depth chain. The depth-limit insufficiency theorem establishes that depth bounds alone cannot prevent this; scope refinement enforcement is independently required.

\textbf{Chain integrity} ensures that the chain topology is stable and tamper-evident. Without it, an adversarial actor could manipulate chain parentage or configuration to redirect accountability or obscure the provenance of a spawned node --- a form of accountability laundering. The immutable parent pointer and node immutability together ensure that every node's lineage and configuration are fixed at provisioning and auditable at any future time.

\textbf{Termination} ensures that the chain cannot grow indefinitely and cannot sustain an infinite spawn-escalate cycle. Without it, an unresolvable exception condition could drive unbounded chain growth exceeding any human's supervisory capacity, defeating the upstream BX3 accountability guarantee by overwhelming the human anchor.

Together these properties guarantee that a Recursive Spawning chain is an accountable, bounded, and auditable refinement process: it grows only in authorized directions (drift prevention), preserves the topology established at provisioning (chain integrity), halts within a \purpose{}-defined depth bound (termination), and maintains a complete forensic record of every decision point (\fact{} layer). This is the operational expression of the upstream BX3 accountability guarantee applied to recursive agent population dynamics.
